% =====================================================================
%  Appendix C --- Azure operations runbook
% =====================================================================
\chapter{Azure Operations Runbook}
\label{app:azure}

This appendix collects the operational procedures of the Microsoft Azure environment hosting the platform. It is intended as a quick reference for the operations team of \hostorg.

\section{VM provisioning}

\subsection{VM specification}
The platform runs on an Azure Linux virtual machine sized for a moderate analytical workload. The VM is provisioned in a Tunisian-or-European region of the Azure cloud, with a managed OS disk and an additional managed data disk dedicated to the PostgreSQL data directory. A static public IP address is associated with the VM.

\subsection{Network Security Group}
The Network Security Group (NSG) attached to the VM applies the following rules:

\begin{itemize}
    \item \textbf{Inbound 22/tcp} (SSH): allowed only from the IP allowlist of the team and of the office network of \hostorg.
    \item \textbf{Inbound 443/tcp} (HTTPS): allowed for the API and the dashboard, behind Nginx with a Let's Encrypt certificate.
    \item \textbf{Inbound 80/tcp} (HTTP): allowed only for the ACME challenge of Let's Encrypt; redirected to HTTPS otherwise.
    \item \textbf{Inbound 5432/tcp} (PostgreSQL): denied. PostgreSQL listens only on \texttt{localhost} and is reachable through SSH tunnels.
\end{itemize}

\subsection{IP allowlist management}
The IP allowlist is maintained as an explicit \texttt{ipset} attached to the NSG. Adding or removing a member of the team is a privileged operation that requires the approval of the platform owner. The history of the allowlist changes is preserved in the operations journal.

\section{SSH access}

\subsection{Account types}
Three account categories share the VM: personal accounts (one per team member, with SSH key authentication only and no shell password), an automation account (\texttt{azureuser}) used by cron and by the systemd services, and the database service account (\texttt{accent\_pipeline}) used internally by the pipeline.

\subsection{SSH tunnel}
The SSH tunnel for PostgreSQL is opened by the team members on their workstation as shown in Listing~\ref{lst:ssh-tunnel-app}. Once the tunnel is open, a notebook or a database client connecting to \texttt{localhost:5432} reaches the analytical PostgreSQL on the VM.

\begin{lstlisting}[style=bash,caption={SSH tunnel for PostgreSQL access.},label={lst:ssh-tunnel-app}]
ssh -i ~/.ssh/azure_id_ed25519 \
    -L 5432:localhost:5432 \
    -N \
    azureuser@<vm-public-ip>
\end{lstlisting}

\subsection{Key rotation}
The SSH keys are rotated annually or upon departure of a team member, whichever comes first. The rotation procedure is recorded in the operations journal.

\section{Pipeline operations}

\subsection{Cron entry}
The cron entry under \texttt{/etc/cron.d/accent-fleet} triggers the incremental pipeline every five minutes (cf.\ Listing~\ref{lst:cron} of Chapter~\ref{chap:deploy}). The output of the run is piped to the system journal under the tag \texttt{accent-fleet}.

\subsection{Inspecting the run log}
The history of the pipeline runs is observed through the SQL query of Listing~\ref{lst:run-log}.

\begin{lstlisting}[style=sql,caption={Inspecting the run log.},label={lst:run-log}]
SELECT run_id, mode, status, started_at, ended_at, rows_processed, error_message
FROM warehouse.etl_run_log
ORDER BY started_at DESC
LIMIT 20;
\end{lstlisting}

\subsection{Triggering a backfill}
A backfill is triggered manually under operator control by Listing~\ref{lst:backfill}. The backfill runs in a \texttt{tmux} session so that the operator can disconnect without interrupting the run.

\begin{lstlisting}[style=bash,caption={Triggering a chunked backfill.},label={lst:backfill}]
tmux new -s backfill
cd /opt/accent-fleet-analytics
.venv/bin/python scripts/run_batch.py --mode backfill --chunk-days 30
# Detach with Ctrl-b, d
\end{lstlisting}

\section{Database operations}

\subsection{Backup}
A nightly logical backup is exported through \texttt{pg\_dump} to \texttt{/opt/backups/}, rotated through a seven-day retention policy. A weekly full snapshot of the data disk is configured at the Azure infrastructure level and retained for four weeks.

\subsection{Credential rotation}
The credentials of the service account \texttt{accent\_pipeline} are rotated quarterly. The rotation procedure updates the password in PostgreSQL, in the \texttt{.env} file of the VM, and restarts the API and the dashboard systemd units to reload the connection pool.

\section{API and dashboard operations}

\subsection{Service management}
The API and the dashboard are managed as systemd units \texttt{accent-api.service} and \texttt{accent-dashboard.service}. Listing~\ref{lst:systemd} shows the canonical commands to inspect, restart and follow the logs of these services.

\begin{lstlisting}[style=bash,caption={Managing the API and dashboard systemd units.},label={lst:systemd}]
sudo systemctl status  accent-api.service
sudo systemctl restart accent-api.service
sudo journalctl -u accent-api.service -f

sudo systemctl status  accent-dashboard.service
sudo systemctl restart accent-dashboard.service
sudo journalctl -u accent-dashboard.service -f
\end{lstlisting}

\subsection{TLS certificate renewal}
The TLS certificate of Nginx is provisioned by Let's Encrypt and is renewed automatically by a daily timer. The renewal status is verified through \texttt{certbot certificates}.

\section{Monitoring and alerting}

\subsection{Pipeline alerts}
A failure of the pipeline (status \texttt{ERROR} in \texttt{etl\_run\_log}) triggers a notification on the operations channel. The notification carries the run identifier, the operating mode and the error message.

\subsection{Data quality alerts}
A failure of any critical check of the validation suite triggers a notification, with a link to the corresponding row of the validation table.

\subsection{Drift alerts}
A statistically significant deviation of the cluster proportions or of the risk band proportions, computed against a one-month baseline, triggers a notification accompanied by a link to the dashboard tile that visualizes the drift.

\cleardoublepage
